% =====================================================================
% SECCIÓN 7: DISCUSIÓN, CONCLUSIONES Y TRABAJO FUTURO
% =====================================================================

% ---------------------------------------------------------------------
% Diapositiva 19: Discusión y Aportes
% ---------------------------------------------------------------------
\section{Conclusiones y trabajo futuro}

\begin{frame}{Discusión integral y aportes de la investigación}
  \framesubtitle{Síntesis de los cuatro hallazgos empíricos y metodológicos principales}

  \begin{columns}[T]
    \column{0.49\textwidth}
    \ugrblockok{1. Superioridad contextual demostrada}{
      BETO supera a los modelos léxicos a nivel de documento ($\Delta F_1 = +0{,}062$, $\Delta\text{AUC} = +0{,}052$, $p < 0{,}0005$). Los 19 rasgos manuales no superan a los n-gramas puros.
    }
    \vskip0.05cm
    \ugrblock{2. Límites de la traducción automática}{
      La traducción neuronal expande la cobertura hacia comunidades de foros ($F_1 = 0{,}934$), pero no transfiere mejoras al dominio nativo de microtexto ($0{,}701$ vs.\ $0{,}717$).
    }

    \column{0.49\textwidth}
    \ugrblock{3. Eficiencia en el balanceo}{
      El submuestreo balanceado 50/50 alcanza un desempeño equivalente a la ponderación de pérdida utilizando un 40\,\% menos de datos, optimizando recursos computacionales.
    }
    \vskip0.05cm
    \ugralert{4. La unidad de análisis importa}{
      En publicaciones individuales aisladas, BETO es ampliamente superior. En perfiles completos con miles de publicaciones, el agregado estadístico equipara el rendimiento (AUC $\ge 0{,}98$).
    }
  \end{columns}

  \note{
    Sintetizamos los cuatro aportes cardinales: demostración de la superioridad de BETO en textos breves; límites empíricos del aumento por traducción; eficiencia del balanceo de clases; y el impacto decisivo de la unidad de análisis entre publicaciones aisladas y perfiles acumulados por usuario.
  }
\end{frame}

% ---------------------------------------------------------------------
% Diapositiva 20: Limitaciones del Estudio
% ---------------------------------------------------------------------
\begin{frame}{Limitaciones de la investigación y consideraciones éticas}
  \framesubtitle{Transparencia metodológica sobre el alcance real del sistema}

  \begin{columns}[T]
    \column{0.49\textwidth}
    \ugralert{Limitaciones metodológicas y técnicas}{
      \begin{itemize}
        \item \textbf{Sin validación clínica individual:} Etiquetas derivadas de redes sociales, no de entrevistas diagnósticas estructuradas (SCID-5 / MINI).
        \item \textbf{Sesgo de autoselección:} Los usuarios activos en redes no representan a la población clínica general.
        \item \textbf{Restricciones de hardware:} GPU local de 8\,GB VRAM limitó el ajuste a 128 tokens y 3 épocas, posponiendo validación cruzada 5-fold en BETO.
      \end{itemize}
    }

    \column{0.49\textwidth}
    \ugrblock{Compromiso ético y privacidad}{
      \begin{itemize}
        \item \textbf{Herramienta no diagnóstica:} Concebida como apoyo al cribado poblacional y vigilancia analítica, nunca como sustituto del criterio médico profesional.
        \item \textbf{Privacidad por diseño:} Anonimización estricta e irreversible mediante hash criptográfico (SHA-256) sin reidentificación de sujetos.
        \item \textbf{Uso responsable:} Delimitación clara del alcance para prevenir decisiones autónomas automatizadas sobre individuos.
      \end{itemize}
    }
  \end{columns}

  \note{
    Asumimos con rigor las limitaciones del trabajo: las etiquetas de redes no equivalen a un diagnóstico clínico formal, y las restricciones de GPU acotaron los modelos probados. Subrayamos que un clasificador de aprendizaje automático solo puede operar como herramienta de apoyo preliminar y nunca debe sustituir el juicio clínico profesional.
  }
\end{frame}

% ---------------------------------------------------------------------
% Diapositiva 21: Trabajo Futuro
% ---------------------------------------------------------------------
\begin{frame}{Líneas de trabajo futuro}
  \framesubtitle{Oportunidades de continuidad y extensión de la investigación}

  \begin{columns}[T]
    \column{0.49\textwidth}
    \ugrblock{1. Modelado y validación clínica}{
      \begin{itemize}
        \item \textbf{Adaptación continua de dominio:} Preentrenamiento no supervisado sobre corpus de salud mental previos al fine-tuning.
        \item \textbf{Validación clínica certificada:} Evaluación sobre cohortes con instrumentos estandarizados (PHQ-9 o DAIC-WOZ).
      \end{itemize}
    }

    \column{0.49\textwidth}
    \ugrblock{2. Temporalidad y explicabilidad}{
      \begin{itemize}
        \item \textbf{Detección temprana secuencial:} Modelado longitudinal (protocolo eRisk) para medir latencia mínima de alertas.
        \item \textbf{Interpretabilidad profunda:} Métodos de atribución (SHAP, Integrated Gradients) y análisis de atención en BETO.
      \end{itemize}
    }
  \end{columns}

  \note{
    Planteamos cuatro líneas de continuidad prioritarias: esquemas de preentrenamiento continuo adaptativo; validación con cohortes clínicas certificadas; modelado secuencial temporal según el estándar de eRisk; y métodos avanzados de explicabilidad en modelos Transformer.
  }
\end{frame}

% ---------------------------------------------------------------------
% Diapositiva 22: Conclusiones Finales
% ---------------------------------------------------------------------
\begin{frame}{Conclusiones finales}
  \framesubtitle{Cierre de la defensa de tesina de grado}

  \ugrblockok{Respuesta a la pregunta central de investigación}{
    \textbf{Es viable y ventajoso} emplear modelos contextuales preentrenados en español (BETO) para detectar indicios de depresión en publicaciones individuales, superando en más de 6 puntos de $F_1$-macro a los métodos clásicos con significancia estadística comprobada ($p < 0{,}0005$).
  }

  \vskip0.1cm
  \begin{itemize}
    \item \textbf{Aporte al estado del arte en español:} Se consolidó un benchmark riguroso, reproducible y abierto con más de 1{,}3 millones de documentos evaluados bajo particiones limpias por usuario.
    \item \textbf{Aporte metodológico crítico:} Se demostró que la traducción automática expande dominios pero no transfiere mejoras al dominio nativo, y que la unidad de análisis es crítica: la ventaja contextual de BETO en textos breves individuales se equipara con modelos lineales al evaluar perfiles de usuario agregados (AUC $\ge 0{,}98$).
    \item \textbf{Replicabilidad garantizada:} Repositorio con código modular, pipelines de preprocesamiento, configuraciones y predicciones disponibles para su verificación.
  \end{itemize}

  \note{
    En conclusión: BETO demuestra una ventaja representacional indiscutible en publicaciones individuales. Sin embargo, la mayor contribución de esta tesina radica en su rigor metodológico: demostrar el impacto de la partición por usuario y cómo la unidad de análisis modifica las conclusiones entre publicaciones individuales y perfiles acumulados.
  }
\end{frame}

% ---------------------------------------------------------------------
% Diapositiva 23: Cierre y Agradecimientos
% ---------------------------------------------------------------------
\begin{frame}[plain]
  \begin{tikzpicture}[remember picture,overlay]
    \fill[ugrAzul] (current page.north west) rectangle (current page.south east);

    \node[anchor=center] at (current page.center) {
      \begin{minipage}{0.85\textwidth}
        \centering
        \includegraphics[height=1.1cm]{logo_ugr.png}\par\medskip
        {\fontsize{22}{26}\selectfont\bfseries\color{white} ¡Muchas gracias por su atención!}\par\medskip
        {\fontsize{12}{15}\selectfont\color{ugrCeleste!40} Preguntas, comentarios y discusión}\par\medskip
        {\color{white}\rule{0.65\textwidth}{0.8pt}}\par\medskip
        {\footnotesize\color{white}
          \textbf{Giuliano Crenna} \quad · \quad \textbf{Bruno Emmanuel Pace}\\[3pt]
          Licenciatura en Ciencia de Datos · Universidad del Gran Rosario (2026)\\[6pt]
          \color{ugrCeleste!30}\texttt{https://github.com/giulicrenna/Deteccion-Temprana-Depresion}
        }
      \end{minipage}
    };
  \end{tikzpicture}

  \note{
    Agradecemos a los evaluadores Jorge Fernández Acuña, Briant Gauna y Alejandro Hernández, a los docentes de la carrera y a la Universidad del Gran Rosario por su constante orientación. Quedamos a total disposición del tribunal para responder a sus preguntas y comentarios. Muchas gracias.
  }
\end{frame}
